
\vspace{0.25cm}
% Dans la partie « Dynamique dans un référentiel non galiléen », l'étude du champ de pesanteur est conduite en supposant le référentiel géocentrique galiléen. De nombreuses applications permettent d'illustrer cette partie : le pendule de Foucault, la déviation vers l'est, les vents géostrophiques, les courants marins ; l'étude statique des marées constitue également une ouverture pertinente.

\begin{tabularx}{\textwidth}[t]{|X|X|}
  \hline 
  \textbf{Notions et contenus} & \textbf{Capacités exigibles} \\
  \hline
  \multicolumn{2}{|l|}{4.2 Actions de contact dans un fluide en mouvement } \\
  \hline 
  Forces de pression. Équivalent volumique. & Exprimer la force de pression exercée par un fluide sur une surface élémentaire. Exprimer l'équivalent volumique des forces de pression à l'aide d'un gradient. \\
  \hline
    Équilibre d'un fluide dans un référentiel non galiléen en translation ou en rotation uniforme autour d'un axe fixe dans un référentiel galiléen. & Établir et utiliser l'expression de la force d'inertie d'entraînement volumique. \\
  \hline 
  Notion d'écoulement parfait et de couche limite. & Exploiter l'absence de forces de viscosité et le caractère isentropique de l'évolution des particules de fluide. Utiliser la condition aux limites sur la composante normale du champ des vitesses. \\
  \hline 
  Relation de Bernoulli pour un écoulement parfait, stationnaire, incompressible et homogène dans le champ de pesanteur uniforme dans un référentiel galiléen. & Établir et utiliser la relation de Bernoulli pour un écoulement parfait, stationnaire, incompressible et homogène dans le champ de pesanteur uniforme dans un référentiel galiléen. \\
  \hline
\end{tabularx}
